\chapter{Background and Related Work}
\label{cha:background}

This chapter reviews the core components referenced in the report and identifies the gaps this project addresses.

\section{PhysTwin and Real2Sim Digital Twins}
\label{sec:bg-phystwin}

PhysTwin \parencite{phystwin} takes synchronised multi-view RGBD video of a controller
interacting with an object, segments the object, and tracks dense motion on both the object and
the controller. It lifts those tracks into 3D with depth and trains a spring--mass \emph{warp}
model in NVIDIA Warp \parencite{warp}. The trained model responds to a controller point in the
simulator and reproduces soft-body behaviour.

The PhysTwin pipeline assumes cameras fixed in space while the object and controller(s) move. This assumption complicates multi-view depth reconstruction.

\section{Monocular and Video Depth Estimation}
\label{sec:bg-depth}

VDA \parencite{vda} is a monocular depth model that produces temporally consistent maps for a
single view. Depth Anything 3 (DA3) \parencite{da3} predicts spatially consistent geometry from an
arbitrary number of images, with or without known camera poses. Multi-view input does not by itself ensure
temporal consistency for a moving scene recorded by stationary cameras. In the configuration tested
here, depth varies between frames. The streaming adaptation processes overlapping chunks, but its
Umeyama alignment becomes singular for larger chunks; duplicating the static views as a dynamic
sequence helps only over a limited window (Chapter~\ref{cha:system-design}).

DUSt3R and MASt3R \parencite{dust3r,mast3r} use pairwise 3D correspondences to recover poses.
Their estimates also vary over time in the tested sequences.

As of the current state of research, no model is both spatially and temporally correct. Models that reason over one frame or several
views flicker between frames, while video-consistency models are smooth but geometrically weaker
and monocular (Chapter~\ref{cha:smoother}).

\section{Camera Pose Estimation}
\label{sec:bg-pose}

Intrinsics and extrinsics serve two purposes here: aligning depth maps with the RGB frames used
for segmentation, and triangulating tracked 2D points into 3D tracks. They are not always
measured during capture.

Pose from the depth model was poor in both respects, misaligning 3D projections and degrading
optimisation and training. Empirically, pairwise correspondences from DUSt3R with a heuristic scale aligned
much better (Chapter~\ref{cha:system-design}).

\section{Temporal Depth Stabilisation}
\label{sec:bg-temporal}

VDA \parencite{vda} adds temporal attention at several decoder scales and introduces a temporal
gradient matching (TGM) loss. That loss supervises frame-to-frame \emph{change} in depth rather
than absolute depth, and inference runs over a segment-wise sliding window.

StableDPT \parencite{stabledpt} attends spatially within each frame in the encoder and
temporally across frames in the decoder. It also adopts the TGM loss.

NVDS and NVDS+ \parencite{nvds,nvdsp} stabilise depth after the fact, using sliding-window
cross-attention together with an optical flow model, but they require metric depth.

\begin{table}[H]
	\caption{Comparison of prior temporal depth stabilisation methods}
	\label{tab:bg-temporal}
	\begin{center}
		\begin{tabular}{@{}>{\raggedright\arraybackslash}p{2.2cm}
		                  >{\raggedright\arraybackslash}p{4.6cm}
		                  >{\raggedright\arraybackslash}p{2.4cm}
		                  >{\raggedright\arraybackslash}p{2.2cm}
		                  >{\raggedright\arraybackslash}p{2.4cm}@{}}
			\toprule
			Method & Temporal mechanism & Supervision & Ground-truth depth & Multi-camera input \\
			\midrule
			VDA & temporal attention at multiple decoder scales,
			                      segment-wise sliding window & TGM loss & yes (lidar) & no (monocular video) \\
			StableDPT & spatial attention per frame, temporal
			                      attention across frames in decoder & TGM loss & yes (lidar) & no (monocular video) \\
			NVDS / NVDS+ & post-hoc stabiliser, sliding-window
			                      cross-attention + optical flow & temporal loss via flow
			                      model & yes (sensor ground truth) & no (monocular video) \\
			Li et al. (2023) & not stated & not stated & not stated & not verified \\
			\bottomrule
		\end{tabular}
	\end{center}
\end{table}


All four methods need ground-truth or metric depth, and all four are monocular: each takes a single
video stream and none accepts synchronised views from several cameras. This project
instead needs a smoother for several synchronised views from an uncalibrated capture, trained
\emph{without ground-truth depth} by treating a video depth model (VDA) as a motion oracle.

\section{Dense Tracking and 3D Reconstruction}
\label{sec:bg-tracking}

CoTracker \parencite{cotracker} tracks points jointly rather than independently, which makes it
more stable than pairwise optical flow. It is the tracker the base pipeline uses, and since it
works in 2D, its tracks are lifted into 3D with the estimated depth. Tracking accuracy is
therefore bounded by depth accuracy.

MVTracker \parencite{mvtracker} tracks across several views, SpatialTrackerV2
\parencite{spatrackerv2} tracks in 3D. On the captures used here, the multi-view tracker performed worse than CoTracker and
the 3D tracker produced visible artefacts (Chapter~\ref{cha:evaluation}).

\label{sec:bg-reconstruction}

SAM and its successors \parencite{sam,sam3,groundedsam} produce promptable, open-vocabulary
masks, and their quality depends strongly on how the object is described. Trellis
\parencite{trellis} was evaluated as a source of mesh and shape priors, while appearance is
reconstructed with Gaussian splatting \parencite{gaussiansplatting}.

\section{Summary of the Gap}
\label{sec:bg-gap}

Each area above is well developed on its own, but few models combine the three properties needed here. No widely used
model is spatially consistent, temporally smooth and multi-camera capable at once. Existing
stabilisers need ground-truth depth and are monocular, and most multi-camera reconstruction
assumes a stationary scene with a moving camera, whereas this project observes a moving object
from static cameras.

DIPhysTwin addresses the gap in two ways. The first is a ground-truth-free smoother that
learns a residual correction to flickering depth, with a video model serving only as a motion
oracle (Chapter~\ref{cha:smoother}). The second is a set of robustness measures that let the
pipeline tolerate imperfect input (Chapter~\ref{cha:system-design}).
